% 本文件由 LaTeX 排版 Agent 根据有效 translation.json 自动生成。
% author=Methodius; work_id=0625; title=选自论复活讲道

\chapter{选自\WorkTitle{论复活}讲道}

% structure: p0001 source=h1 target=omitted reason=page-title-already-rendered-by-parent

% structure: p0002 source=h2 target=section reason=relative-heading-rank; base=h2

\section{第一部分}

一、天主并未造恶，祂也绝非恶的作者；但凡受造的，既蒙赐予\OriginalTerm{自由意志}{free-will}以守护祂公义所定的律法，却未能持守这律法，便称为恶。违抗天主，僭越那与自由意志相符的正义之界限，正是最严重的过犯。\par

二、如今这个问题已被提出并解答了：那\Quote{\OriginalTerm{皮衣}{coats of skins}}\BibleRef{创 3:21}并非身体。然而，我们仍要再说一次，因为仅提一次是不够的。在预备这些皮衣之前，原祖自己承认他有骨有肉；因为他看见女人被领到他面前时，便喊道：\BibleQuote{这才是}\BibleRef{创 2:23}\Quote{骨中骨，肉中肉}。又说：她应称为\Person{女人}，因为是从男人取出的。为此，人要离开父母，与妻子结合，二人成为一体。因我无法容忍某些人的妄言，他们无耻地强解圣经，好使他们那\Quote{复活没有肉体}的见解得到支持；他们妄称有\OriginalTerm{理性的骨肉}{rational bones and flesh}，并用各种\OriginalTerm{寓意解经}{allegorizing}之法翻来覆去地改变经文。基督也证实了这些事当按字面理解：当法利塞人问及休妻之事时，祂回答：\BibleQuote{你们没有读过吗？那起初造人的，造了他们一男一女，且说：为此，人要离开父亲}\BibleRef{玛 19:4}等等。\par

三、然而，认为在永恒状态中身体不会与灵魂共存，因为身体是\OriginalTerm{捆绑与锁链}{bond and fetters}，这种想法显然是荒谬的；依他们之见，我们这些将在\OriginalTerm{光明王国}{kingdom of light}中生活的人，就不该永远被判作腐化的奴仆。既然这问题已充分解决，他们那将肉体定义为灵魂锁链的论断已被驳斥，那么\Quote{肉体不会复活，免得我们若重新穿上它，便在光明王国里成为囚犯}这一论点也就瓦解了。\par

四、为使人类不致成为不死或永存的恶——倘若罪在人内掌权，如同在不朽的身体中滋生，并有不朽的养料，那么情况就会如此——天主因此宣告他必死，并以\OriginalTerm{可朽性}{mortality}给他披上。因为这就是\OriginalTerm{皮衣}{coats of skins}所指的：藉着身体的解体，罪得以从根部完全消灭，甚至不留下一丝根须，以免罪孽的新芽再次萌发。\par

五、正如一棵无花果树，生长在宏伟圣殿的建筑中，长得高大，枝繁叶茂的根盘绕着石块的所有接缝，它不断生长，直等那石头从它生长的地方被松开，它才被完全拔除；因为无花果树被拔走后，石头仍可放回原位，圣殿得以保存，不再支撑那导致自身毁灭的原因；而无花果树连根拔起，便枯死了；同样，天主的圣殿——人，当他滋养了罪，如同一棵\OriginalTerm{野无花果树}{wild fig-tree}时，建造者天主便适时地用死亡来遏止，按圣经的话：\BibleQuote{杀死}\BibleRef{申 32:39}\Quote{并使人活}，为叫罪恶枯干死亡之后，肉体能像修复的圣殿，以同样的部分再次复活，完好无损且不朽，而罪则被彻底消灭。因为当身体还活着、尚未经历死亡时，罪也必与它同活，罪的根隐藏在我们内，即使外表上因训诫与警告的打击而受抑制，它仍存在；否则，我们在领洗之后就不会犯罪了，因为我们本该完全彻底无罪。但如今，即使在信德之后，在接触圣化之水的时候之后，我们也常常陷于罪恶。因为无人能夸口自己如此无罪，甚至没有恶念。因此，罪如今因信德被抑制并平息，以致不结出有害的果实，然而并未连根拔除。目前，我们克制它的嫩芽，如恶念，\BibleQuote{恐怕有\OriginalTerm{苦根}{root of bitterness}长起来扰乱}\BibleRef{希 12:15}我们，不让它的叶子张开并抽出枝条；而天主圣言如斧头，砍断下面生长的根。但将来，连恶念本身也要消失。\par

VI. 然而，既然在此类事上需要许多例证，现在让我们特别从这一观点来探究，不达更有力的解释和明证，就不停止。看来，好像一位卓越的工匠亲自用黄金或其他材料铸成一座高贵的像，各部分比例优美，而他突然发现那像被某个卑鄙之人毁坏了——那人因嫉妒不能容忍像的美丽，便糟蹋了它，以此空尝嫉妒之乐。最明智的阿格拉奥丰，请注意：如果工匠希望他付出如此多心血和劳力的作品完全不受损害，他必会将其熔化，恢复原状。但若他不重新铸造，也不重建，而任其存留，加以修补和复原，那么那像经过火与锻造，便不能再保持不变，而会走样和损耗。因此，若他想要它完美无瑕，就必须打碎重铸，以便因背叛和嫉妒所加的一切毁损，通过打碎重铸得以消除，而像复归完好、纯净，与先前相同，尽可能与自身相似。因为正在原匠手中的像，即使再度熔化，也不可能丧失，因它可以复原；但瑕疵和损伤则可以摆脱，因它们熔化而无法复原；因为在每件艺术品中，最优的工匠所寻求的不是瑕疵或失败，而是其作品的匀称与正确。在我看来，天主的计划与我们中通行的相同。因祂看见祂最美好的作品——人，被嫉妒的背叛所败坏，祂以对人的慈爱，不能容忍他留在那般境况，免得他永远有瑕疵，永负罪责；而是将他再次分解为原始材料，以便通过重塑，他内的一切瑕疵得以消逝不见。因为前例中像的熔化相当于后例中身体的死亡和解体，前例中材料的重塑相当于后例中死后的复活；正如耶肋米亚先知也这样说，他向\Emph{犹太人}讲话说：\BibleQuote{我下到陶工的家，看，他正在轮子上工作。他手中的器皿坏了，他又照他眼中的好另做一器皿。上主的话传给我说：以色列家！我难道不能像这陶工一样对待你们吗？看，你们在我手中，就像陶泥在陶工手中一样。}\BibleRef{耶 18:3}\par

VII. 我要提醒你注意：如我所说，在人犯罪后，那只伟大的手并不满足于将自己的作品——被恶者出于嫉妒邪恶地损伤了的——作为胜利的标记留下；而是将它浸湿、化解为泥土，如同陶工打碎器皿，以便通过重塑，其中一切瑕疵和伤痕得以消失，并可重新成为完美中悦的器皿。\par

VIII. 但若说宇宙将被完全毁灭，海、空气和天空不复存在，这并不令人满意。因为整个世界将被天上的烈火淹没，经过焚烧以达到净化和更新；但它不会完全毁灭和腐朽。因为若世界不存在比存在更好，那么天主创造世界时为何选取更糟的途径？但天主并未枉然行事，也未做最糟的事。因此，天主为存在和延续而安排了受造界，正如\WorkTitle{智慧篇}所证实的，说：\BibleQuote{因为天主造了万物，使它们存在；世界的世代是健康的，它们内没有毁灭的毒药。}\BibleRef{智 1:14}保禄也清楚见证，说：\BibleQuote{受造之物热切地等待天主子民的显示。因为受造之物被屈服于空虚，并非出于自愿，而是由于那使之屈服者的希望；因为受造之物本身也要从腐朽的奴役中获得释放，进入天主子女光荣的自由。}\BibleRef{罗 8:19}他说受造之物被屈服于空虚，并期望它从这种奴役中获得释放，因为他意指此世界为受造之物。因为不是那看不见的，而是那看得见的，才属于腐朽。因此，受造之物在更新为更美、更适宜的状态后，存留下去，在复活时因天主的子女而欢喜雀跃；它如今为他们的缘故呻吟劳苦，也等着我们从身体的腐朽中得救赎，就如经上所载：\BibleQuote{耶路撒冷啊！抖下尘土，起来，坐下吧！}\BibleRef{依 52:2}以及当我们起来、摆脱肉身的必死性、脱离罪恶之后，它也要脱离腐朽，不再屈服于空虚，而屈服于正义。依撒意亚也说：\BibleQuote{正如我所造的新天新地，在我面前存留，上主说，你们的后裔和名字也必这样存留。}\BibleRef{依 66:22}又说：\BibleQuote{那创造诸天、铺张大地、且使之存立的上主这样说：祂并未将地创造为空虚，而是塑造它为人居住。}\BibleRef{依 45:18}实际上，天主并非枉然或无目的地设立宇宙以致毁灭，如那些思想软弱者所说，而是为存在、为居住、为延续。因此，在万物受到焚烧和摇撼之后，大地和天空必须再次存在。\par

IX. 但如果我们的反对者说，既然宇宙并未被毁灭，主为何却说\BibleQuote{天地要过去}？\BibleRef{玛 24:35}又先知说\BibleQuote{诸天必如烟云消散，大地必要像衣服一样破旧}？\BibleRef{依 51:6}我们回答：因为圣经惯常将世界从现状转变到更美好、更光荣状态的过程称为\Quote{毁灭}，因为万物在变化为更辉煌的状态时，其先前形态便消失了；因为圣经中并无矛盾或荒谬之处。因为经上说，不是\Quote{世界}而是\Quote{这世界的局面}要过去，\BibleRef{格前 7:31}所以圣经惯常将从较早的形态转变为更美好、更适宜的状态称为\Quote{毁灭}；正如人们将孩子形态转变为完人称为毁灭，因为孩童的身量转变为成人的大小和美丽。我们可以预期受造物将在烈焰中逝去，仿佛毁灭，以便得以更新；但并非彻底消灭，好使我们这些被更新的人居住在一个更新的世界中，不再尝受忧愁；正如经上说：\Quote{你发出你的气息，它们便受造，你便使地面更新。}此后天主为那包围它的东西提供适宜的温度。因为在这现今的时代之后，大地仍将存在，必定有它的居民，他们不再有死亡的危险，也不娶不嫁，不生子女，而是像天使一样，在一切幸福中生活，没有变化或朽坏。因此，如果不再有空气、大地或其他任何东西，那么讨论我们身体将如何生活是愚蠢的。\par

X. 此外，除了已说的，还有一点值得考虑，因为它非常误导人——如果我们可以就如此重要的事直言不讳的话，\Person{阿格拉奥丰}——因为你说主曾清楚宣告，那些将获得复活的人，那时将如同天使。\BibleRef{玛 22:30}你提出这个反对：天使没有肉身，因此享有最大的幸福和光荣。那么，既然我们要变得与天使同等，我们就必须像他们一样脱去肉身，成为天使。但你忽略了这一点，我优秀的朋友：那从无中创造并安排宇宙的他，将不朽者的本性不仅分派给天使和服役者，也分派给率领者、宝座者和掌权者。因为天使的族类是一样，率领者和掌权者的族类是另一样；因为不朽者并非全都属于同一等级、构造、部族和家族，而是有族类和部族的区别。革鲁宾不会脱离自己的本性而采取天使的形状；同样，天使也不会采取其他形状。因为他们不能成为别的，只能成为他们之所是及所被造的模样。此外，人既按万物的原始秩序被任命居住世界，并统治其中的一切，当他成为不朽时，也绝不会从人改变成天使或其他任何形状；因为天使也不会从原形改变成别的形状。因为基督来临的时候，并没有宣告人性在不朽时应被重塑或转变成另一种本性，而是宣告应恢复至堕落前的状态。因为每一个受造物必须留在自己适当的位置，使任何事物都不缺少，而一切得以充满：天堂充满天使，宝座充满掌权者，光体充满服役者；还有更神圣的处所，无玷无瑕的光体，与革鲁宾一同侍奉最高议会，支撑宇宙；以及人的世界。因为如果我们承认人会改变成天使，那么我们也必须说天使会改变成掌权者，而这些又会变成这变成那，直到我们的论证过远而陷入危险。\par

XI. 天主并没有好像造人很差劲或在塑造他时犯了错误，而后决定造一个天使，如同最差劲的工匠后悔自己的工作那样；祂也没有在原先想造天使却失败之后，再造人；因为这将显示软弱，等等。那么，如果祂希望人成为天使而不是人，为什么祂造人而不是造天使呢？是因为祂不能吗？这样想就是亵渎。或者祂忙于制造较差的，以致在更好的事情上怠惰？这也是荒谬的。因为祂在创造善事上不会失败，也不会拖延，更不会无能；相反，祂有能力按自己的意愿和时机行事，因为祂本身就是能力。因此，正因为祂意欲人是人，所以起初造人如此。但如果祂如此意欲——既然祂意欲的是善——那么人是善的。现在，人据说由灵魂和身体组成；那么他不能没有身体而存在，只能有身体，除非在人类之外产生了另一种人。因为所有不朽者的品级都必须由天主保存，而人也在其中。\BibleBook{}{智慧篇}说：\BibleQuote{天主造人，为使其不朽，并使他成为自己永恒性的肖像。}\BibleRef{智 2:23}所以身体不会朽坏；因为人是由灵魂和身体组成的。\par

十二、因此，请注意，这些正是主愿意教训那些不相信肉身复活的\Person{撒杜塞人}的事。因为这就是\Person{撒杜塞人}的意见。因此，他们编造了关于那妇人和七个弟兄的比喻，试图对肉身的复活提出质疑，经上说：\BibleQuote{有来到祂跟前}，\BibleRef{玛 22:23}据说，\BibleQuote{撒杜塞人，他们说没有复活}。\newline{} 那么，基督若没有肉身的复活，而只有灵魂得救，祂就应当赞同他们的意见，认为那是正确而卓越的。但事实并非如此，祂回答说：\BibleQuote{在复活的时候，他们也不娶也不嫁，而是像天上的天使一样}，\BibleRef{玛 22:23}这并不是因为他们没有肉身，而是因为他们不再嫁娶，从此不朽。祂论到我们在这一点上近似天使，就是说，如同天上的天使一样，我们在乐园中也如此，不再花时间于婚宴或其他庆典，而是在基督的指引下，瞻仰天主并培育生命。因为祂并没有说\Quote{他们要成为天使}，而是像天使一样，例如，正如经上所载，被戴上了荣耀和尊贵为冠冕；他们与天使稍有差异，却又近乎天使。正如祂若说：观察天空的美丽秩序，夜晚的寂静，以及被月亮的天空之光所照亮的一切，\Quote{月亮如同太阳一样发光}。那么，我们不应当说祂断言月亮就是太阳，而是像太阳。同样，那不是金子，却接近金子的性质，被称为不是金子，而是像金子。但如果是金子，就会被说是金子，而不是像金子。既然它不是金子，却接近它的性质，并有它的外表，就被说成是像金子。同样，当祂说圣人们将在复活时像天使一样，我们并不认为祂断言他们那时实际上将成为天使，而是接近天使的状况。因此，最不合理的是有人说：\Quote{既然基督宣告圣人们在复活时显现像天使，所以他们的身体不会复活}，尽管所使用的词语本身清楚地证明了实际情况。因为\Quote{复活}一词并不适用于那没有跌倒的，而是适用于那跌倒而又起来的；正如先知所说：\BibleQuote{我也要重建\Person{达味}已倒塌的帐幕}。\BibleRef{亚 9:11} 如今，这备受渴望的灵魂的帐幕已经倒塌，沉入\BibleQuote{地里的尘土}之中。\BibleRef{达 12:2} 因为所安放的，不是那没有死的，而是那死了的。但死的是肉身；灵魂是不死的。那么，如果灵魂不死，而身体是尸体，那些说有复活，却没有肉身复活的人，就否认了任何复活；因为起来的是那已经跌倒安放的，而不是那仍然站着的；正如经上所载：\BibleQuote{那跌倒的岂不再起来？那转离的岂不归回？}\BibleRef{耶 8:4}\par

十三. 既然肉身被造于\OriginalTerm{不朽}{incorruption}与\OriginalTerm{朽坏}{corruption}的交界处，本身既不是前者也不是后者，却因快乐而被朽坏胜过了，尽管它是不朽的作品和财产；因此它变成了可朽坏的，并被置于地上的尘土中。然而，当它被朽坏所胜，因背命而被交付于死亡时，天主并没有将它留给朽坏作为产业被战胜；而是借着\OriginalTerm{复活}{resurrection}战胜死亡后，重新将它交付于不朽，为的是让朽坏不能获得不朽的产业，而是不朽获得朽坏的产业。因此宗徒如此回答说：\BibleQuote{因为这必朽坏的必须穿上不朽，这必死的必须穿上不死。} \BibleRef{格前 15:53}现在这必朽坏的和必死的穿上\OriginalTerm{不死}{immortality}，除了是那\BibleQuote{播种在朽坏中，复活在不朽中，} \BibleRef{格前 15:42}——因为灵魂并非可朽坏或必死的；但这必死和朽坏的是属于肉身的——为的是\BibleQuote{我们怎样带着属土的肖像，也怎样带着属天的肖像。} \BibleRef{格前 15:49}因为我们所带的\OriginalTerm{属土肖像}{earthy image}正是这个：\BibleQuote{你本是尘土，仍要归于尘土。} \BibleRef{创 3:19}但\OriginalTerm{属天肖像}{heavenly image}是从死者中的复活和不朽，为的是\BibleQuote{基督怎样借着父的光荣从死者中复活，我们也怎样在新生中行走。} \BibleRef{罗 6:4}但如果有人认为属土肖像就是肉身本身，而属天肖像却是肉身之外的另一个属神身体；就让他首先考虑：基督，那属天的人，显现时，带着和我们相同的肢体形式和相同的肉身肖像，借此，那本不是人的祂成了人，为的是\BibleQuote{正如在亚当内众人都死了，照样在基督内众人也都要复活。} \BibleRef{格前 15:22}因为如果祂取了肉身，是为了其他缘故，而不是为了释放肉身并复活它，那为什么祂要多余地取了肉身，如同祂既无意拯救它，也不愿复活它？但天主之子不做任何多余的事。祂取了仆人的形状，并非无用，而是为了复活并拯救它。因为祂真的成了人，真的死了，不是仅仅表面如此，而是为真正显示祂是\OriginalTerm{从死者中复活的首生者}{first begotten from the dead}，将属土的变为属天的，将必死的变为不死的。因此，当保禄说\BibleQuote{血肉不能承受天主的国，} \BibleRef{格前 15:50}他并非对肉身的重生给出贬低的意见，而是要教导：天主的国，即永生，并非由身体拥有，而是身体由生命拥有。因为如果天主的国，即生命，被身体拥有，那么生命就会被朽坏耗尽。但现在生命拥有那正在死去的，为的是\BibleQuote{死亡被生命吞灭以致胜利} \BibleRef{格前 15:54}，而可朽坏的将被视为不朽和不死的产业，同时它变得不受束缚，并自由于死亡和罪，却成为不死的奴隶和仆人；为使身体成为不朽的产业，而不是不朽成为身体的产业。\par

十四、那么，若人从如此渺小、先前毫无存在、实际上只是潮湿、蜷缩且微不足道的一滴液体——其实是从虚无——被带入存在，那么，人从先前已存在的人中重新进入存在，岂不更加容易？因为，使曾经存在后又腐朽的事物重新被造，并非像从虚无中造出从未存在之物那样困难。现在，如果我们愿意展示从男人排出的精液，并将其与一具尸体并排放置，两者各置一处，旁观者看这两样摆着的东西，会认为哪一样更可能变成一个人——那完全虚无的液滴，还是那已有形状、大小和形质的尸体？因为，那本是虚无之物，仅因天主愿意，便成为人；那么，那已存在且被塑造圆满之物，若天主愿意，岂不更应再成为人？神学家梅瑟在《肋未纪》中以神秘意义引入帐棚节，其目的是什么？难道是为了让我们向天主守节，如同那些以低俗眼光理解圣经的犹太人那样？好像天主喜欢那些用果实、树枝和叶子装饰、随即枯萎凋零的帐棚。我们不能这样说。那么请告诉我，帐棚节的目的何在？它是为了指向我们这真实的帐棚——它因违犯法律而堕入腐朽，被罪恶拆毁——天主应许要重新组合它，并在不朽中使之复活，好使我们真正在复活时，为光荣祂而庆祝那伟大而著名的帐棚节；那时，我们的帐棚将按不朽与和谐的完美秩序被组合，从尘土中不朽地复活；那时，枯骨，\BibleRef{则 37:4}，按最真实的预言，将听到声音，并由天主——造物主和完美工匠——接合，祂将更新肉体并重新缚上，不再用起初维系它的那种系带，而是用永远不腐、不可解散的系带。因为我曾在利西亚的奥林帕斯山看到，山顶上有火自发地从地中喷出；火旁长着一棵杨梅树，如此茂盛、翠绿、荫蔽，以至于人们可能以为是不竭的水流滋养了它的生长，而非实际情况。因此，虽然事物的本性质易朽，其身体会被火吞噬，且一旦具有可燃本性之物就不可能不受火影响，但这棵树非但未被烧毁，反而比平时更加茁壮翠绿，尽管它本性可燃，且火焰就在其根旁燃烧。我确实从邻近的树林中折了些树枝抛到火喷出的地方，它们立刻着火，烧成了灰。那么，请告诉我，那连太阳的热都经受不住、若不浇水就会枯萎的树，为何被如此炽热的火焰包围却不被烧毁，反而生机勃勃？这奇事的含义是什么？天主将此作为将要来临之日的预兆和引子，使我们更确切地知道，当万物被天火淹没时，那些以贞洁和正义为特征的体魄，将被祂取走，不受火的任何伤害，如同凉水一般。因为，仁慈而慷慨的主啊，\Quote{那受造的万物，事奉祢——造物主，为惩罚不义者而增强力量，为信任祢者的利益而减弱力量；}\BibleRef{智 16:24}按祢的旨意，火冷却，不伤害祢决定保存的任何事物；同样，水比火更猛烈地燃烧，没有任何事物能抵抗祢不可征服的大能与权柄。因为祢从虚无中创造了万物；因此，祢也随己意改变和转化万物，因为它们都属于祢，而祢是唯一的天主。\par

十五、宗徒确实在将栽种和浇灌归于技艺、土地和水之后，唯独将生长归于天主，他说：\Quote{可见，栽种的不算什么，浇灌的也不算什么，只在那使之生长的天主。}\BibleRef{格前 3:7}因为他知道，那作为天主首生者、万有的父母和工匠的智慧，将万物带入世界；古人称她为本性和眷顾，因为她以不断的安排和照顾，赐予万物诞生与生长。\Quote{因为，}天主的智慧说，\Quote{我父至今工作，我也工作。}\BibleRef{若 5:17}正因此，撒罗满称智慧为万有的工匠，因为天主在任何方面都不匮乏，而是能丰盛地创造、制作、变化并增加万物。\par

十六、天主，创造了万物、并眷顾和照料万物者，从地上的尘土取来，造了我们外在的人。\par

% structure: p0019 source=h2 target=section reason=relative-heading-rank; base=h2

\section{第二部分}

\Emph{论复活第二讲。}\par

例如，我们君王的肖像，即使不是用更珍贵的材料——金或银——制成，也为众人所尊敬。因为人们并不会因尊重那些用昂贵材料制成的肖像，而轻视那些用次等材料制成的，而是尊敬世上每一肖像，即使是白垩或青铜的。无论谁对其中任何一者出言不逊，都不会被宣告无罪，仿佛他只是冒犯了黏土，也不会被定罪为轻视了金子，而是因对君主和主本身无礼。天主天使的肖像，用金子铸成，即元首和掌权者的肖像，我们是为祂的光荣和荣耀而制作的。\par

% structure: p0022 source=h2 target=section reason=relative-heading-rank; base=h2

\section{第三部分}

\Emph{一、摘自论复活讲道。}\par

I. 请阅读殉道主教\Person{圣美多德}的\WorkTitle{论复活}一书，以下为其节选。身体不是灵魂的桎梏，正如\Person{奥力振}所认为的那样，灵魂也并非因在身体内而被先知\Person{耶肋米亚}称为\Quote{束缚}。因为他定下原则：身体不妨碍灵魂的活动，反而，身体随着灵魂运行，并在灵魂所托付的一切事上与灵魂合作。但我们应如何理解\Person{额我略神学家}和许多其他人的见解呢？\par

II. \Person{奥力振}说，身体在堕落之后被给予灵魂作为桎梏，而此前灵魂没有身体；但我们所穿戴的这个身体是我们犯罪的原因；因此他也称它为桎梏，因为它会阻碍灵魂行善。\par

III. 如果说身体在堕落之后被给予灵魂作为桎梏，那么它必然是被给予作为对于邪恶或对于良善的桎梏。但它不可能作为对于良善的桎梏；因为没有医生或工匠会给出了差错的东西一个导致更大错误的补救，更何况天主呢。因此，它只能是对于邪恶的桎梏。但我们分明看到，起初，\Person{加音}穿着这个身体杀了人；很明显，其后继者陷入了何等的邪恶。所以，身体不是对于邪恶的桎梏，也根本不是桎梏；灵魂也并非在堕落之后才第一次穿上身体。\par

IV. 就本性而言，人最真实地既不能被说成是没有身体的灵魂，也不能被说成是没有灵魂的身体，而是由灵魂与身体结合成一体之美的存在。但\Person{奥力振}说唯独灵魂是人，\Person{柏拉图}也是如此。\par

V. 人与其他生物之间是有区别的；它们被赋予了多种自然形式和形状，如同有形可见的自然力量在天主的命令之下所产生的数量；而人却被赋予了天主的模样和肖像，每一个部分都精确地完成，按照圣父和独生圣子的原始相似。现在我们应考虑圣人如何说明这一点。\par

VI. 他说，雕刻家\Person{菲迪亚斯}在制作了象牙的皮萨神像之后，吩咐在它面前倒油，以便尽可能保护它不朽。\par

VII. 他说，正如\Person{雅典纳哥拉}也曾说过的，魔鬼是一个受造于天主之旁的精神体，靠近物质，当然与其他天使一样；他被托付管理物质和物质的形态。因为，按照天使最初的构成，他们是天主在其眷顾中所造，为照料宇宙；为使天主对万物施行完美而普遍的监督，并保持对万有的至高权能和权力（因为万有依赖祂存在），而被委派的天使负责具体事务。其余的天使留在天主创造并安置他们的位置上；但魔鬼骄横，对我们起了嫉妒，在所托付的职分中行为邪恶；后来那些迷恋肉欲并与人的女儿们发生非法交合的天使也是如此。因为对他们，如同对人一样，天主也赐予了他们自己的选择。这应如何理解？\par

VIII. 他说，皮衣象征死亡。因为他说到\Person{亚当}，当全能的天主看到因背叛，他——一个不朽的存在——变成了邪恶，正如他的欺骗者魔鬼一样，天主因此准备了皮衣；这样，当他如此，可以说，穿上必朽性时，他里面一切邪恶就可在身体的分解中死去。\par

IX. 他认为\Person{圣保禄}有两次启示。因为宗徒，他说，并不像那些懂得语言精妙的人所认为的，将乐园置于第三层天上，当他这样说：\BibleQuote{我知道这样一个人被提到第三层天上；我也知道这样一个人——或在身内，或在身外，天主知道——他被提到乐园里。}\BibleRef{格后 12:2} 这里他指出他看见了两次启示，明显被提了两次，一次到第三层天，一次到乐园。因为\BibleQuote{我知道这样一个人被提到}这话肯定了他自己得到了关于第三层天的启示；而随后的话：\BibleQuote{我也知道这样一个人——或在身内，或在身外，天主知道——他被提到乐园里}，表明关于乐园的另一个启示给了他。他之所以这样说是由于他的对手从宗徒的话中断定乐园仅仅是一个概念，因它位于天之上，以便得出乐园中的生命是无形的结论。\par

X. 他说，行恶或避免行恶在于我们的能力；否则我们就不会因行恶而受罚，也不会因行善而受赏；但邪恶思想的有无并不取决于我们自己。因此，就连\Person{圣保禄}也说：\BibleQuote{我所愿意的，我偏不作；我所憎恨的，我反而去作。}\BibleRef{罗 7:15} 这就是说，\Quote{我的思想不是我所愿意的，而是我所不愿意的。} 现在他说，想象邪恶的习惯因肉体死亡的临近而被根除——因为正是为此，天主为罪人设定了死亡，使邪恶不致永远存留。\par

但这一陈述是什么意思？值得注意的是，我们其他教父也曾这样说过。\Emph{其含义何在？}既然那些面对死亡的人在那一刻既没有罪的增加也没有罪的减少？\par

\Emph{II. 同一论著中某些宗徒话语的纲要。}\par

1. 请从同一篇讲道中读一下某些宗徒话语的简明解释。那么，让我们看看我们试图对宗徒所说的话是什么意思。因为他的这句话：\BibleQuote{我以前没有法律是活着的，}\BibleRef{罗 7:9}指的是在律法之前，我们的原祖父母在乐园里所过的生活，不是没有身体，而是有身体，正如我们上面所表明的；因为我们活着没有\OriginalTerm{私欲偏情}{concupiscence}，完全不知道它的攻击。因为没有一条我们应当遵循的法律，也没有能力决定我们应当采取何种生活方式，从而可以被公正地赞许或指责，这被认为可以免除一个人的指控。因为一个人不会\OriginalTerm{贪欲}{lust}那些没有禁止他的东西，即使他贪欲了，也不会被责备。因为贪欲不是指向那些在我们面前、受我们支配的东西，而是指向那些在我们面前、却不在我们能力范围之内的东西。因为一个人怎么会关心一件既不禁忌也不必要的东西呢？因此经上说：\BibleQuote{我若不认识贪欲，除非法律说过：\InnerQuote{你不可贪恋。}}\BibleRef{罗 7:7}因为当（我们的原祖父母）听到：\BibleQuote{分别善恶树上的果子，你不可吃，因为你吃的日子必定死，}\BibleRef{创 2:17}他们就产生了贪欲，并摘取了它。所以经上说：我若不认识贪欲，除非法律说过：你不可贪恋；他们也不会想吃，除非曾说过：\BibleQuote{你不可吃它。}因为从那以后，罪就趁着机会欺骗了我。因为当法律赐下时，魔鬼就有能力在我里面激起贪欲；\BibleQuote{因为没有法律，罪是死的；}\BibleRef{罗 7:8}这意味着\Quote{当法律没有赐下时，罪不能犯。}但我是在法律之前活着、无可指责的，没有一条必须遵行的诫命；\BibleQuote{但诫命来到，罪又活了，我就死了。那本来叫人活的诫命，反倒叫我死。}\BibleRef{罗 7:9}因为在天主赐下法律、命令我当作何事、不当作何事之后，魔鬼就在我里面激起了贪欲。因为天主赐给我的应许，这是为了生命和不朽，这样我若顺从它，就可以拥有常青的生命和喜乐，直到不朽；但若违背它，就会导致死亡。但是魔鬼——他称之为罪，因为他是罪的作者——趁着机会，藉着诫命欺骗我违命，欺骗并杀害了我，从而使我要受那谴责：\BibleQuote{你吃的那一天，必定要死。}\BibleRef{创 2:17}\BibleQuote{这样看来，法律是圣洁的，诫命也是圣洁、公义、良善的；}\BibleRef{罗 7:12}因为它赐下，不是为了伤害，而是为了安全；因为我们不要以为天主造了任何无用或有害的东西。那么怎么样呢？\BibleQuote{那良善的是叫我死吗？}\BibleRef{罗 7:13}即那作为律法赐下的，是为了成为最大善的原因吗？\BibleQuote{断乎不是！}\BibleRef{罗 7:13}因为使我屈从于朽坏的原因不是天主的法律，而是魔鬼；好叫那藉着善行恶的显明出来，使那恶的发明者成为并显明为众罪人中最伟大的。\BibleQuote{因为我们知道法律是\OriginalTerm{属神的}{spiritual}；}\BibleRef{罗 7:14}所以它在任何方面都不能伤害任何人；因为属神的事远远脱离无理性的贪欲和罪。\BibleQuote{但我是\OriginalTerm{属血肉的}{carnal}，是已经卖给罪了。}\BibleRef{罗 7:14}意思是：但我既是属血肉的，被置于善恶之间作为一个\OriginalTerm{自愿主体}{voluntary agent}，就能够凭我的意愿选择。因为\Quote{看，我将生命和死亡摆在你面前；}意思是：违命属神的律法，即诫命，会导致死亡；顺从属血肉的律法，即蛇的计谋，也会如此；因为藉着这样的选择，\Quote{我被卖给}了魔鬼，落在罪之下。因此，邪恶仿佛围攻我，依附我，住在我里面，正义将我交付给那恶者，因为我违犯了律法。所以那些话：\BibleQuote{我所做的，我自己不明白；我所恨恶的，我倒去做。}\BibleRef{罗 7:15}并不理解为行恶，而只是思想恶。因为思想或不思想不当之事，不在我们能力之内，但行动或不行动于我们的思想，则在能力之内。因为我们不能阻止思想进入我们的心智，因为当它们从外面被灌输进来时，我们便接受了；但我们能够克制自己不顺从它们、不付诸行动。因此，我们有能力愿意不想这些事情；但不能使它们消散，不再进入头脑；因为这不在我们的能力之内，正如我所说；这就是那句话的意思：\BibleQuote{我所愿意的善，我反不作；}\BibleRef{罗 7:19}因为我不愿意思想那些伤害我的事；因为这个善完全是无可指责的。但\BibleQuote{我所愿意的善，我反不作；我所不愿意的恶，我倒去作；}\BibleRef{罗 7:19}不愿意思想，却思想了我不愿意的事。试想，达味不正是为这些事而恳求天主，悲叹自己思想了那些他不愿意的事吗？\Quote{求祢洗涤我，使我洁净，清除我的隐恶。求祢不让骄傲人辖制我，我便完全，免犯大罪。}宗徒在另一处也说：\BibleQuote{攻破那些阻拦人认识天主的自高之事，又将人所有的心意夺回，使它顺服基督。}\BibleRef{格后 10:5}\par

第二，若有人胆敢反对这一陈述，回答说：宗徒教导我们，我们不仅憎恨思想上的恶，而且我们做我们所不愿做的，并且在做的同时憎恨它，因为他说：\BibleQuote{我所愿意的善，我不去做；我所不愿意的恶，我倒去做。}\BibleRef{罗 7:19}如果这样说的人说的是真理，那么我们就请他解释：宗徒所憎恨且不愿做却做了的恶是什么？他所愿做却不做的善是什么？反过来，是否他每次愿行善时，他都不做他所愿的善，而做他所不愿的恶？并且他如何能在劝勉我们摆脱一切罪恶时说：\BibleQuote{你们当效法我，如同我效法基督一样？}\BibleRef{格前 11:1}因此，他所指的是那些他已经提到、他并不愿去做、只是思想的那些事，而不是实际去行的。否则，他如何能完全效法基督呢？因此，如果我们没有那些反对我们、与我们争战的人，那将是极好且最愉快的；但既然这是不可能的，我们就不能做我们所愿的。因为我们不愿有那些引诱我们陷于情欲的人，那样我们就能不费力地得救；但所愿的并没有成就，反而我们所不愿的却发生了。因为如我所说，我们必须经受考验。所以，我的灵魂啊，我们不要向那恶者屈服；而要穿上天主的全副武装，即我们的护盾，让我们有\BibleQuote{正义的甲冑，并以和平福音的预备作为鞋穿在脚上；此外，还要拿起信德的盾牌，使你们能扑灭那恶者的一切火箭；并戴上救恩的头盔，带上圣神的剑，即天主的圣言，}使你们能抵抗魔鬼的诡计；\BibleQuote{推翻一切对傲慢自高的认识，并摧毁一切反对基督知识的高墙，}\BibleRef{格后 10:5}\BibleQuote{因为我们作战不是对抗血与肉；}\BibleRef{弗 6:12}\BibleQuote{我所做的，我不认可；我所愿意的，我倒不做；我所憎恨的，我反而去做。若我做我所不愿做的，我就同意法律是善的。那么现在不再是做那事的我，而是住在我内的罪。因为我知道在我内——即在我肉身内——没有善。}\BibleRef{罗 7:15}这话说的很对。因为请回想我们已表明的：从人离正道、违背法律的时候起，罪因他的悖逆而生，就住在他内。因为这样，骚动就被激起，我们充满了激荡和异样的想象，被掏空了天主的灵感，充满了狡猾的蛇注入我们内的肉欲。因此，天主为了我们的缘故发明了死亡，为毁灭罪，免得罪在我们这些不死者中继续存在，如我所说，成为不死的。当宗徒说：\BibleQuote{因为我知道在我内——即在我肉身内——没有善}时，他以此指明罪因违命藉着贪欲住在我们内；由此，如同嫩芽，快乐的想象在我们周围出现。因为我们内有两种思想：一种来自肉身内的贪欲，如我所说，来自恶灵的诡计；另一种来自符合诫命的法律，它被植入我们内作为自然法律，当我们按我们的心思喜爱天主的法律时，它激发我们的思想向善，因为这就是内在的人；另一种则来自魔鬼的法律，按照住在肉身内的贪欲。因为那与天主的法律作战并对抗的，即是反对心意向善的倾向，正是那激起肉身的、感官的冲动去悖逆者。\par

三、因为在我看来，宗徒在这里清楚提出了三条法律：一条符合植入我们内的善，他明确称之为\OriginalTerm{理智的法律}{law of the mind}。一条来自恶的攻袭，常使灵魂倾向于欲念的幻想，他说，这法律与理智的法律交战。\BibleRef{罗 7:23} 第三条，符合罪，因私欲而居于肉身中，他称之为\Quote{居于我们肢体内的\OriginalTerm{罪恶的法律}{law of sin}}；\BibleRef{罗 7:23} \OriginalTerm{恶者}{Evil One}催促它，常常激起它攻击我们，驱使我们走向不义和恶行。因为在我们自身内，似乎有一种更好的和一种更坏的。当那本性上更好的将要胜过更坏的时，整个心灵便被引向善；但当更坏的增强并压倒时，人反而被催迫趋向邪恶的想象。因此，宗徒祈求从中获得解救，视之为死亡和毁灭；先知也同样说：\Quote{求你洗净我隐秘的过犯。}同样的话也由以下经文表明：\Quote{因为我按\OriginalTerm{内在的人}{inward man}喜悦天主的法律，但我看出在我肢体内另有一条法律，与我理智的法律交战，并把我掳去，隶属于那在我肢体内的罪恶的法律。我真不幸！谁能救我脱离这\OriginalTerm{死亡的身体}{body of this death}呢？}\BibleRef{罗 7:22} 他这样说，并非指身体是死亡，而是指那居于他肢体内的罪恶的法律，因违命而隐藏在我们内，不断欺骗灵魂，使之趋向不义的死亡。他立刻加上，清楚表明他渴望脱离何种死亡，以及谁解救了他：\Quote{我感谢天主，借着耶稣基督。}\BibleRef{罗 7:25} 并且应当思考，倘若他像你所假设的那样，\Person{阿格拉风}，称这身体是死亡，他就不会随后提到基督解救他脱离如此巨大的邪恶。因为在那种情况下，基督的来临会给我们带来何等奇怪的事呢？宗徒怎能说，他能借着基督的来临脱离死亡？因为基督来到世界之前，死亡是所有人的定命。所以，\Person{阿格拉风}，他并非说这身体是死亡，而是说那因私欲而居于身体内的罪，天主借着基督的来临已解救他脱离了这罪。\Quote{因为生命之圣神的法律在基督耶稣内使我脱离了罪恶与死亡的法律；}以致\Quote{那使耶稣从死者中复活者，也要借着那居于你们内的圣神，使你们有死的身体复活；}祂\Quote{定了罪}那在身体内的罪，使之灭绝；\Quote{使那法律的正义}——这使我们趋向善的本性法律，且符合诫命——得以点燃并彰显。因为本性法律所不能行的善事，\Quote{因它软弱无力，}被身体内的私欲所战胜，天主却赐予力量完成，\Quote{派遣了自己的儿子，带着\OriginalTerm{有罪肉身的形状}{likeness of sinful flesh}；}以致罪被定罪而灭绝，使其在肉身中永不结果实，好使本性法律的正义得以满全，充盈于那些不随从肉欲、而随从圣神的私欲和引导者的顺命中。\Quote{因为生命之圣神的法律}（即福音），不同于先前的法律，借着宣讲引导我们顺服和罪过的赦免，使我们脱离了罪恶与死亡的法律，完全战胜了统治我们肉身的罪。\par

四、他说，植物并非从土中得到滋养或增长。因为他说，让任何人都想想，土如何能改变并被吸收到树木的质料中。因为若是这样，那么原来周围的土通过树根被吸入树木的整个范围，树木生长的地方，必然会被挖空；所以他们关于物体流动的说法是荒谬的。因为土怎能先通过根进入植物的树干，然后经过它们的管道进入所有枝条，变成叶子和果实呢？有很多大树，如雪松、松树、冷杉，每年结很多叶子和果实；人们可以看到它们并没有消耗周围的土来增加树的体积和材质。因为如果真是土通过根上升变成木头，那么它们周围土所在的地方必然会被挖空；因为干燥的物质不像湿润的物质那样能够流入并填补移动所留下的空间。此外，还有无花果树和其他类似的植物，常常生长在纪念建筑上，但它们从未将整个建筑吸收到自己体内。但如果有人多年来收集它们的果实和叶子，他会发现它们的体积已变得比纪念建筑上的土大得多。因此，认为土被消耗成果实和叶子的收成是荒谬的；即使它们都是由土所造，也仅仅是将土作为其根基和位置。因为面包的制造离不开磨、地方、时间和火；但面包并非由这些东西中的任何一种制成。其他成千上万的事也是如此。\par

五、现在，\Person{奥力振}的追随者们引出这段经文：\BibleQuote{因为我们知道，如果我们这地上的帐棚式的房屋被拆毁，}\BibleRef{格后 5:1}等等，藉以否认身体的复活，声称\Quote{\OriginalTerm{帐棚}{tabernacle}}是身体，而\Quote{\OriginalTerm{非人手所造的房屋}{house not made with hands}}\Quote{在天上}是我们的\OriginalTerm{属神的衣服}{spiritual clothing}。因此，\Person{圣美多迪乌}说，这\OriginalTerm{地上的房屋}{earthly house}必须隐喻地理解为我们在世短暂的存在，而非这帐棚；因为如果你决定将身体视为被拆毁的地上的房屋，请告诉我们那房屋被拆毁的帐棚是什么？因为帐棚是一回事，帐棚的房屋是另一回事，而我们这些拥有帐棚的又是另一回事。他说：\Quote{因为，如果我们这帐棚的地上的房屋被拆毁，}\BibleRef{格后 5:1}——借此他指出，灵魂是我们自己，身体是帐棚，而帐棚的房屋隐喻地象征今生肉体的享乐。那么，如果身体的今生如房屋般被拆毁，我们将拥有那在天上非人手所造的。他说：\Quote{非人手所造的}，以指出区别；因为可以说今生是手造的，因为今生的一切职业和追求都是由人的手进行的。因为身体，作为天主的化工，不可说是手造的，因为它不是由人的技艺形成的。但如果他们说它是手造的，因为它是天主的化工，那么我们的灵魂，以及天使和天上的属神的衣服，也都是手造的；因为所有这些也都是天主的化工。那么，什么是手造的房屋？正如我所说的，是由人手维持的短暂存在。因为天主说：\BibleQuote{你必汗流满面，才得糊口；}\BibleRef{创 3:19}当那生命被拆毁时，我们便拥有那非人手所造的生命。正如主所指示的，祂说：\BibleQuote{要藉着那不义的钱财结交朋友，使到你们匮乏的时候，他们可以接你们到永久的帐幕里。}\BibleRef{路 16:9}因为主在那里所称的\Quote{帐幕}，宗徒在这里称为\Quote{衣服}。祂在那里所称的\Quote{不义的朋友}，宗徒在这里称为\Quote{被拆毁的房屋}。因此，当我们今生的日子终结时，那些我们在这不义的生命和这\Quote{世界}（它\Quote{卧在那恶者手下}，\BibleRef{若一 5:19}）中所获得的善行，必将接引我们的灵魂；同样，当这必朽的生命被拆毁时，我们将拥有复活之前的住处——也就是说，我们的灵魂将与天主同在，直到我们领受那为我们预备的、永不倒塌的新房屋。因此我们也\Quote{叹息}\Quote{并不是愿意脱下}身体，\Quote{乃是愿意穿上}\BibleRef{格后 5:4}来世的生命。因为我们渴望\Quote{穿上}的\Quote{天上的房屋}是不朽性；当我们穿上它时，一切软弱和必死性都将被它完全\Quote{吞灭}，被永无穷尽的生命所消耗。\BibleQuote{因为我们行事为人是凭着信德，不是凭着眼见；}\BibleRef{格后 5:7}这就是说，我们仍在凭信德前行，以昏暗的悟性观看那超越的事物，而非清晰地看待，以便我们能看见、享有并置身其中。\BibleQuote{弟兄们，我告诉你们说，血肉之身不能承受天主的国，朽坏也不能承受不朽坏。}\BibleRef{格前 15:50}他所说的血肉，不是指血肉本身，而是指灵魂中朝向淫欲之乐的非理性冲动。因此，当他说\Quote{血肉之身不能承受天主的国}时，他加以解释：\Quote{朽坏也不能承受不朽坏。}现在，\OriginalTerm{朽坏}{corruption}不是那被朽坏的东西，而是那朽坏的东西。因为当死亡得胜时，身体便陷入朽坏；但当生命仍存于其中时，它便保持\OriginalTerm{不朽坏}{incorruption}。因此，既然血肉是朽坏与不朽坏之间的边界，既不是朽坏也不是不朽坏，它便因快乐而被朽坏所战胜，尽管它是不朽坏的作品和产业。因此它成了朽坏的奴隶。那么，当它被朽坏征服并交给死亡作为惩罚时，祂没有任凭它被战胜并作为遗产交给朽坏；而是再次藉着复活战胜死亡，使它恢复不朽坏，好使朽坏不能继承不朽坏，而是不朽坏继承那朽坏的。因此宗徒回答说：\BibleQuote{这必朽坏的必须穿上不朽坏，这必死的必须穿上不死。}\BibleRef{格前 15:53}但必朽坏的和必死的穿上不朽坏和不死，这不是别的，正是那在朽坏中播种的，在不朽坏中复活？\BibleRef{格前 15:42}因为\BibleQuote{我们既然带有属土的肖像，也将带有属天的肖像。}我们曾带有的\Quote{\OriginalTerm{属土的肖像}{image of the earthly}}指的是这句话：\Quote{你本是尘土，仍要归于尘土。}而\Quote{\OriginalTerm{属天的肖像}{image of the heavenly}}则是从死者中复活和不朽坏。\par

VI. 纳波利的犹斯定，一位既未远离宗徒时代，也未远离宗徒德行的人，论道：那属朽坏的被继承，而生命却继承；血肉死去，而天国活着。因此，当保禄说\BibleQuote{血肉不能承受天国}时，他并非似乎在轻视血肉的重生，而是教导：天主的国，即永生，并非由身体所承受，而是身体被生命所承受。因为，倘若天主的国，即生命，被身体所承受，那么生命就会被腐朽所吞没。但如今，生命继承那属朽坏的，为叫死亡被生命吞没而获胜利，并使那属可朽的成为不朽的财产；它从死亡和罪恶中获得自由，成为不朽的奴仆和臣属，好叫身体成为不朽的财产，而非不朽成为身体的财产。\par

VII. 对于经文\BibleQuote{在基督内死了的人要先复活；然后我们这些活着的人}，圣美多迪乌斯如此解释：那些是我们的身体；因为灵魂就是我们自己，我们复活，从地上取回那已死的；如此，我们与他们一同被提到云中，迎接主，便可荣耀地庆祝复活的灿烂节庆，因为我们领受了永久的帐幕，不再死亡，也不再解体。\par

VIII. 他说，我看见在奥林帕斯（奥林帕斯是吕基亚的一座山）上，有一团火从山顶的地上自发燃起；火旁有\OriginalTerm{普拉格诺斯}{Puragnos}植物，如此茂盛、翠绿、多荫，好似它从泉水生出。何故？虽它们本性属可朽的，且其身体为火所焚，这植物却不仅未被烧毁，反而更加茂盛，尽管它本质上易燃，且火正在其根部燃烧？然后我从周围的树林里取树枝，抛入火焰喷涌之处，它们即刻爆燃，化为灰烬。这矛盾有何含义？天主将此指定为那将要来临之日的一个标记和序幕，为叫我们知道：当万物被火淹没时，那拥有贞洁与正义的身体将穿越火焰，如同穿过冷水。\par

IX. 他说，请考虑：有福的若望说\BibleQuote{海交出其中的死者；死亡和阴府交出其中的死者}，这岂不意味着各元素交出各部分以重建每人？海指水元素；阴府指空气，源自\OriginalTerm{无形}{\GreekText{ἀειδές}}，因其不可见，正如奥利振所言；死亡指地，因为死者被安葬于其中；故此在圣咏中，它也被称为\BibleQuote{死亡的尘土}，基督说祂被带入\BibleQuote{死亡的尘土}。\par

X. 他说，因为凡是纯粹由空气与纯粹由火所构成，且与天使存有本质相同的，不可能有土与水的本性；否则，它就属于土性了。奥利振已指明，人的身体——那将要复活的身体——就是这样的本性，由这样的元素构成；他还说这身体将是属神的。\par

XI. 他问：复活的身体将有何形貌？若按照他的说法，这人的形貌作为无用之物将完全消失；而它本是活物中最可爱的形貌，因为天主自己就采用这个形貌，正如最明智的保禄所解释：\BibleQuote{男人不可蒙头，因为他是天主的肖像和光荣}；与之一致，天使的理性身体也以此形貌排列？它将是圆形、多边形、立方体还是金字塔形？因为形貌种类繁多；但这是不可能的。那么，我们该如何看待那声称：神肖的形貌要被摒弃为更卑贱的？因为他自己承认灵魂像身体，且人复活时没有手脚？\par

XII. 他说，转化即是恢复为不受苦、有荣耀的状态。因为如今身体是欲望和卑贱的身体，故此达尼尔被称为\BibleQuote{欲望之人}。但那时，它将被改变为不受苦的身体，不是通过肢体的排列变化，而是通过不再贪求肉体的享乐。\par

然后他驳斥奥利振说：奥利振因此认为，同样的血肉不会归还给灵魂，而是每个人的形貌——根据现在血肉所藉以区分的样式——会印在另一个属神的身体上复活；如此，每个人将再次显现同一形貌；这就是所应许的复活。因为，他说，物质身体是流动的，绝不停留在自身之内，而是围绕着那区分其形貌的样式（形状由此样式所包含）逐渐耗损和替换；因此复活只能是形貌的复活。\par

十三．过了一会儿，他又说：\Person{奥力振}，如果你主张肉身的复活只是外表上转变为属神的身体，并引\Person{梅瑟}和\Person{厄里亚}的显现作为最强有力的证据，说他们离世后显现，保持了与生前毫无二致的外貌，那么所有人的复活也必如此。但\Person{梅瑟}和\Person{厄里亚}是在基督受难及复活之前，以你说的这种形相复活并显现的。那么，先知和宗徒们又怎能称基督为\Quote{死者中的首生者}呢？因为如果相信基督是死者中的首生者，祂就是最先复活的。但\Person{梅瑟}在基督受难前就已向宗徒显现，具有你所说复活所成全的形相。因此，没有肉身的形相复活是不成立的。要么如你所教导的，有形相的复活，那么基督就不再是\Quote{死者中的首生者}，因为灵魂在死后已先于祂以这种形相显现；要么祂确实是首生者，那么就绝不可能有人在祂之前被认为配得复活而不至再死。但如果没有人先于祂复活，而\Person{梅瑟}和\Person{厄里亚}向宗徒显现时并没有肉身，只有肉身的表象，那么肉身的复活便清楚显明了。因为仅仅以形相来设定复活是极其荒谬的，既然灵魂离肉身之后，从未脱去那形相——而你声称这形相会复活。但如果那形相与灵魂同在，如同\Person{梅瑟}和\Person{厄里亚}的灵魂那样，既不会如你所想的消灭，也不会毁灭，而是处处与它们同在；那么，这从未堕落过的形相就绝不能说是复活了。\par

十四．如果有人觉得这不可接受，反问说：如果在基督降入阴府之前没有人复活，那么为何记载有数人在祂之前复活？其中有\Place{漆冬}寡妇的儿子、\Place{叔能}妇人的儿子，以及\Person{拉匝禄}。我们必须说：这些人复活是为了再死；但我们所谈论的是那些复活后永不再死的人。如果有人对\Person{厄里亚}的灵魂存疑，说圣经记载他是带着肉身被提升的，而我们又说他是脱去肉身向宗徒显现的，那么我们必须说，承认他带着肉身向宗徒显现，反而更有利于我们的论点。因为这一事例表明，肉身是可以承受不朽的，正如\Person{哈诺客}被提升所证明的。因为如果肉身不能承受不朽，它就不可能长时间保持无知觉的状态。那么，如果他带着肉身显现，那便是在他死后真正发生的，但绝非是从死者中复活。我们这样说，是同意\Person{奥力振}所说的：同样的形相在死后赋予灵魂；当灵魂与肉身分离时，这是最不可能的事，因为肉身的形相早已因变化而败坏，正如熔化的雕像在完全解体之前，其形相也已败坏。因为性质不能脱离质料而独立存在；围绕铜像的形相在熔化时与雕像分离，就不再具有实体性的存在。\par

十五．既然形相在死亡中与肉身分离，那么让我们思考一下，所谓分离有多少种方式。一件事物与另一件事物的分离，可以是在动作和实体上，也可以是在思想上；或者在动作上，但不在实体上。例如，若有人将混合在一起的小麦和大麦分开：就它们分离的运动而言，称为动作上的分离；就它们分离开后各自独立而言，称为实体上的分离。当我们把质料与其性质分开，或把性质与质料分开时，便是思想上的分离。当一件事物与另一件事物分离后不再存在，没有实体性时，便是动作上的分离而非实体上的分离。在机械学中也可观察到类似情况，例如当人观看一尊雕像或一匹熔化的铜马时，会看到其自然形相正在变化，变成另一种形状，原有的形相消失。因为若有人将塑造成人形或马形的作品熔化，他会发现形相的外表消失，而质料本身仍存留。因此，主张形相将丝毫不受损坏地复活，而印有形相的身体却被毁灭，是站不住脚的。\par

十六．但他却说，情况正是如此；因为身体将在属神的身体中改变。因此，必须承认，原来的形相并不会如初地复活，因为它已经与肉身一同改变和败坏。即使它改变成属神的身体，那也不是原本的实体，而是在以太般的身体中塑造的某种相似物。然而，如果复活的既不是同一个形相，也不是同一个身体，那么便是以另一个替代了第一个。因为相似之物既然不同于它所相似之物，就不能等同于它所仿效的那个原初之物。\par

十七．此外，他说：那外貌或形相，正是以形相的特有特征显示出肢体的同一性。\par

十八、当\Person{奥利振}将\Person{厄则克耳}先知论死者复活的话作寓意解释，并曲解为以色列子民从\Place{巴比伦}被掳中归回时，圣人在反驳他时，说了诸多话后，也说了这些话：他们既未获得完全的自主，也未以更大的能力战胜敌人而再居\Place{耶路撒冷}；当他们屡次有意建造（圣殿）时，却被别的民族阻止。因此，他们几乎用了四十六年才建造起来，而\Person{撒罗满}从根基起只用七年就完成了。但在这事上何需多言？因为从\Person{拿步高}及以后统治\Place{巴比伦}的人时起，直到\Place{波斯}远征\Place{亚述}、\Person{亚历山大}的帝国、以及\Place{罗马}人向犹太人挑起的战争，\Place{耶路撒冷}被敌人攻陷了六次。这事由\Person{若瑟夫}记载，他说：\Quote{\Place{耶路撒冷}在\Person{维斯帕先}统治的第二年被攻取。以前它已被攻取过五次，但现在是第二次被摧毁。因为\Place{埃及}王\Person{阿索柴乌斯}，之后有\Person{安提约古}，接着\Person{庞培}，然后\Person{索修斯}与\Person{黑落德}攻取了城并焚烧了它；但在他们之前，\Place{巴比伦}王征服并毁灭了它。}\par

十九、他说\Person{奥利振}持有他所反驳的这些意见。关于\Person{拉匝禄}和富翁，可能存在疑问。较为简单的人以为这些话是说两人都在肉身内为活着时所做的事受到报应；但较为精确的人认为，既然复活后没有人留在生命中，这些事不会在复活时发生。因为富翁说：我有五个兄弟……免得他们也来到这个痛苦的地方，请派\Person{拉匝禄}去告诉他们这里的事。因此，如果我们追问关于\Quote{舌头}、\Quote{手指}、\Quote{\Person{亚巴郎}的怀抱}以及在那里斜躺的事，或许灵魂在变化中会得到一个与它粗糙属地的肉身相似的形状。那么，如果任何一位已眠者被记载曾显现，同样地，他显现时是带着他在肉身时的形状。此外，当\Person{撒慕尔}显现时，显然他显现时穿着肉身；如果我们被证明灵魂的本质是无形体、且其自身显示自身的论证所迫，这一点尤其必须承认。但在救主显现之前，在世界终结之前，因而也在复活之前，那在痛苦中的富翁和那在\Person{亚巴郎}怀抱中受安慰的穷人，据说一个在阴府受罚，另一个在\Person{亚巴郎}的怀抱中受安慰；这教导说，现在已经在变化中，灵魂复活成为身体。因此，圣人说如下：\Person{奥利振}主张灵魂在离开此处后具有一个与这个可感觉的身体相似的形状，那么他是否像\Person{柏拉图}一样主张灵魂是无形体的？那在离开世界后据说需要一个运载工具和一件外衣，以免被找到赤裸的东西，它本身怎么可能不是无形体的呢？但如果它是无形体的，它不是也必须无苦难吗？因为它属于无形体，它也是无苦难和不动摇的。那么，如果它不为任何非理性的欲望所困扰，它也不为痛苦或受苦难的身体所改变。因为无形体者不能与身体共感，身体也不能与无形体者共感，如果灵魂如所说的那样看似无形体的话。但如果它与身体共感，正如显现者的证据所证明的，它就不能是无形体的。因此，唯有天主被颂扬为无始、独立、不疲倦的本性；祂是无形体的，因此是不可见的，因为\BibleQuote{从来没有人见过天主}。但灵魂作为理性的形体，被万物的创造者和父安排成理性可见的肢体，接受了这一印记。因此，在阴府中，如同\Person{拉匝禄}和富翁的情形，他们被说成有舌头、手指和其他肢体；这并不是说他们另外有一个不可见的身体，而是说灵魂本身，当完全脱去其外罩时，按其本质就是那样的。\par

二十、圣人末了说：\BibleQuote{因为基督死而复生，正是为作死人和活人的主，}这话必须理解为指灵魂和身体；灵魂是\Emph{活着的}，因为是不朽的；身体是\Emph{死了的}。\par

二十一、既然人的身体比其他活物更为尊贵，因为据说它是天主的双手所造，且成为了理性灵魂的载体；那么它为何如此短命，甚至比一些无理性的受造物还短命？难道不是显然它的长寿存在将在复活之后吗？\par

\SourceRecord{Translated by William R. Clark.；Ante-Nicene Fathers；Vol. 6.；Edited by Alexander Roberts, James Donaldson, and A. Cleveland Coxe.；Buffalo, NY: Christian Literature Publishing Co.,；1886.；<http://www.newadvent.org/fathers/0625.htm>.}
